\documentclass[10pt, twocolumn, landscape, a4paper]{article}

\usepackage[margin=1in]{geometry}
\usepackage[utf8]{inputenc}
\usepackage[spanish,es-noshorthands]{babel}
\usepackage{amssymb, amsmath, amsbsy}
\usepackage{mathpazo}
\usepackage{tikz}
% \usetikzlibrary{angles,quotes}
\usepackage[pdftex]{hyperref}
\hypersetup{pdftitle={Triángulo en coordenadas},pdfauthor={Luis Carrillo Gutiérrez}}

\fontfamily{cmss}\selectfont % sin mathpazo
\renewcommand{\familydefault}{\sfdefault}
\pagestyle{empty} % no número de página

\begin{document}

% Regular polygon in  cartesian coordinate system w/vertex

\begin{center}
	\begin{tikzpicture}
		% 18 - 16 = 2
		\coordinate (A) at (2, 2.01);
		% 16.26 - 16 = 0.26
		\coordinate (B) at (.26, -1.01);
		% 19.74 - 16 = 3.74 
		\coordinate (C) at (3.74, -1.01);
		\draw (A) -- (B) -- (C) -- cycle;
		% \draw (0, 0) -- (A);
	\end{tikzpicture}
\end{center}
\end{document}